\section{Non-integer powers}

\subsection{Definition}

\begin{definition}
\label{def.fps.HasConstantTermOne}
\lean{AlgebraicCombinatorics.FPS.HasConstantTermOne}
\leanhelper
\leanok
An FPS $f\in K[[x]]$ has \emph{constant term one} if $[x^0]f=1$.
We write $f\in K[[x]]_1$ for this condition.
\end{definition}

\begin{lemma}
\label{lem.fps.hasConstantTermOne-mul}
\lean{AlgebraicCombinatorics.FPS.hasConstantTermOne_mul}
\leanhelper
\leanok
If $f,g\in K[[x]]_1$, then $fg\in K[[x]]_1$.
\end{lemma}

\begin{proof}
\leanok
$[x^0](fg) = ([x^0]f)([x^0]g) = 1\cdot 1 = 1$.
\end{proof}

\begin{lemma}
\label{lem.fps.hasConstantTermOne-one-add-X}
\lean{AlgebraicCombinatorics.FPS.hasConstantTermOne_one_add_X}
\leanhelper
\leanok
$1+x\in K[[x]]_1$.
\end{lemma}

\begin{proof}
\leanok
$[x^0](1+x) = 1$.
\end{proof}

\subsubsection{The Log map}

\begin{definition}
\label{def.fps.logSeries}
\lean{AlgebraicCombinatorics.FPS.logSeries}
\leanhelper
\leanok
The \emph{logarithm series} (or $\overline{\log}$) is the FPS
\[
\overline{\log} := \sum_{n\geq 1} \frac{(-1)^{n-1}}{n} x^n
= x - \frac{x^2}{2} + \frac{x^3}{3} - \cdots
\in K[[x]].
\]
Its constant term is~$0$.
\end{definition}

\begin{lemma}
\label{lem.fps.constantCoeff-logSeries}
\lean{AlgebraicCombinatorics.FPS.constantCoeff_logSeries}
\leanhelper
\leanok
$[x^0]\overline{\log} = 0$.
\end{lemma}

\begin{proof}
\leanok
Immediate from the definition.
\end{proof}

\begin{definition}
\label{def.fps.fpsLog}
\lean{AlgebraicCombinatorics.FPS.fpsLog}
\leanhelper
\leanok
For $f\in K[[x]]_1$, we define
\[
\operatorname{Log}(f) := \overline{\log} \circ (f-1).
\]
Since $f$ has constant term~$1$, the FPS $f-1$ has constant term~$0$,
so the substitution is well-defined.
\end{definition}

\begin{lemma}
\label{lem.fps.fpsLog-one}
\lean{AlgebraicCombinatorics.FPS.fpsLog_one}
\leanhelper
\leanok
$\operatorname{Log}(1) = 0$.
\end{lemma}

\begin{proof}
\leanok
$\operatorname{Log}(1) = \overline{\log}\circ (1-1) = \overline{\log}\circ 0 = 0$,
since every term in the substitution sum is zero.
\end{proof}

\begin{lemma}
\label{lem.fps.constantCoeff-fpsLog}
\lean{AlgebraicCombinatorics.FPS.constantCoeff_fpsLog}
\leanhelper
\leanok
For $f\in K[[x]]_1$, we have $[x^0]\operatorname{Log}(f)=0$.
\end{lemma}

\begin{proof}
\leanok
The constant term of a substitution $g\circ h$ where both $[x^0]g=0$ and $[x^0]h=0$
is~$0$.
\end{proof}

\begin{theorem}
\label{thm.fps.fpsLog-mul}
\lean{AlgebraicCombinatorics.FPS.fpsLog_mul}
\leanhelper
\leanok
For $f,g\in K[[x]]_1$,
\[
\operatorname{Log}(fg) = \operatorname{Log}(f) + \operatorname{Log}(g).
\]
\end{theorem}

\begin{proof}
\leanok
Let $h = \operatorname{Log}(fg) - \operatorname{Log}(f) - \operatorname{Log}(g)$.
We show $h'=0$ and $[x^0]h = 0$ using the derivative characterization:
$(d/dx)\operatorname{Log}(f) = (d/dx \overline{\log})\circ(f-1)\cdot f'$,
and the key identity $(d/dx\,\overline{\log})\circ(f-1)\cdot f = 1$
(since $d/dx\,\overline{\log} = 1/(1+x)$ and substituting $f-1$ gives $1/f$).
The Leibniz rule for $(fg)'$ then shows $h'=0$.
Since $[x^0]h = 0-0-0 = 0$ and $h'=0$, uniqueness gives $h=0$.
\end{proof}

\subsubsection{The Exp map}

\begin{definition}
\label{def.fps.fpsExp}
\lean{AlgebraicCombinatorics.FPS.fpsExp}
\leanhelper
\leanok
For $g\in K[[x]]_0$ (FPS with constant term~$0$), we define
\[
\operatorname{Exp}(g) := \exp\circ\, g
= \left(\sum_{n\in\mathbb{N}} \frac{x^n}{n!}\right)\circ g.
\]
Since $g$ has constant term~$0$, the substitution is well-defined.
\end{definition}

\begin{lemma}
\label{lem.fps.fpsExp-zero}
\lean{AlgebraicCombinatorics.FPS.fpsExp_zero}
\leanhelper
\leanok
$\operatorname{Exp}(0)=1$.
\end{lemma}

\begin{proof}
\leanok
Only the $m=0$ term in the substitution sum is nonzero, giving~$1$.
\end{proof}

\begin{lemma}
\label{lem.fps.hasConstantTermOne-fpsExp}
\lean{AlgebraicCombinatorics.FPS.hasConstantTermOne_fpsExp}
\leanhelper
\leanok
For $g\in K[[x]]_0$, we have $\operatorname{Exp}(g)\in K[[x]]_1$.
\end{lemma}

\begin{proof}
\leanok
The constant term of $\operatorname{Exp}(g) = \exp\circ g$ is $[x^0]\exp = 1/0! = 1$,
since all higher powers of $g$ contribute~$0$ to the constant term.
\end{proof}

\begin{theorem}
\label{thm.fps.fpsExp-add}
\lean{AlgebraicCombinatorics.FPS.fpsExp_add}
\leanhelper
\leanok
For $u,v\in K[[x]]_0$,
\[
\operatorname{Exp}(u+v) = \operatorname{Exp}(u)\cdot\operatorname{Exp}(v).
\]
\end{theorem}

\begin{proof}
\leanok
Follows from the multiplicativity of the exponential map, established
in Mathlib's Exp/Log theory for formal power series.
\end{proof}

\subsubsection{The power definition}

\begin{definition}
\label{def.fps.power-c}
\lean{AlgebraicCombinatorics.FPS.fpsPow}
\leantarget
\leanok
Assume that $K$ is a commutative $\mathbb{Q}$-algebra.
Let $f\in K\left[  \left[  x\right]  \right]  _{1}$ and $c\in K$. Then, we
define an FPS%
\[
f^{c}:=\operatorname{Exp}\left(  c\operatorname{Log}f\right)  \in K\left[
\left[  x\right]  \right]  _{1}.
\]
\end{definition}

\begin{lemma}
\label{lem.fps.fpsPow-zero}
\lean{AlgebraicCombinatorics.FPS.fpsPow_zero}
\leanhelper
\leanok
For any $f\in K[[x]]$, we have $f^0 = 1$.
\end{lemma}

\begin{proof}
\leanok
$f^0 = \operatorname{Exp}(0\cdot\operatorname{Log}f) = \operatorname{Exp}(0) = 1$.
\end{proof}

\begin{lemma}
\label{lem.fps.hasConstantTermOne-fpsPow}
\lean{AlgebraicCombinatorics.FPS.hasConstantTermOne_fpsPow}
\leanhelper
\leanok
For $f\in K[[x]]_1$ and $c\in K$, we have $f^c\in K[[x]]_1$.
\end{lemma}

\begin{proof}
\leanok
Since $[x^0]\operatorname{Log}f = 0$, we have $[x^0](c\operatorname{Log}f) = 0$,
so $\operatorname{Exp}(c\operatorname{Log}f)\in K[[x]]_1$.
\end{proof}

\begin{lemma}
\label{lem.fps.fpsPow-one}
\lean{AlgebraicCombinatorics.FPS.fpsPow_one}
\leanhelper
\leanok
For $f\in K[[x]]_1$, we have $f^1 = f$.
\end{lemma}

\begin{proof}
\leanok
$f^1 = \operatorname{Exp}(1\cdot\operatorname{Log}f) = \operatorname{Exp}(\operatorname{Log}f) = f$,
using the inverse property $\operatorname{Exp}\circ\operatorname{Log} = \operatorname{id}$
on $K[[x]]_1$.
The formal proof bridges the local definitions of $\operatorname{Log}$ and $\operatorname{Exp}$
to Mathlib's Exp/Log framework
(via Lemmas~\ref{lem.fps.fpsLog-eq-Log-val} and~\ref{lem.fps.fpsExp-eq-Exp-val})
and then applies the inverse property.
\end{proof}

\subsection{Rules of exponents}

\begin{theorem}
\label{thm.fps.power-c.rules}
\lean{AlgebraicCombinatorics.FPS.fpsPow_add, AlgebraicCombinatorics.FPS.fpsPow_mul, AlgebraicCombinatorics.FPS.fpsPow_pow}
\leantarget
\leanok
Assume that $K$ is a commutative $\mathbb{Q}%
$-algebra. For any $a,b\in K$ and $f,g\in K\left[  \left[  x\right]  \right]
_{1}$, we have
\[
f^{a+b}=f^{a}f^{b},\ \ \ \ \ \ \ \ \ \ \left(  fg\right)  ^{a}=f^{a}%
g^{a},\ \ \ \ \ \ \ \ \ \ \left(  f^{a}\right)  ^{b}=f^{ab}.
\]
\end{theorem}

\begin{proof}
\textbf{First rule ($f^{a+b}=f^af^b$):}
\begin{align*}
f^{a+b} &= \operatorname{Exp}((a+b)\operatorname{Log}f) \\
&= \operatorname{Exp}(a\operatorname{Log}f + b\operatorname{Log}f) \\
&= \operatorname{Exp}(a\operatorname{Log}f)\cdot\operatorname{Exp}(b\operatorname{Log}f)
\quad\text{(by }\operatorname{Exp}(x+y)=\operatorname{Exp}(x)\operatorname{Exp}(y)\text{)} \\
&= f^a \cdot f^b.
\end{align*}

\textbf{Second rule ($(fg)^a = f^a g^a$):}
\begin{align*}
(fg)^a &= \operatorname{Exp}(a\operatorname{Log}(fg)) \\
&= \operatorname{Exp}(a(\operatorname{Log}f + \operatorname{Log}g))
\quad\text{(by }\operatorname{Log}(fg)=\operatorname{Log}f+\operatorname{Log}g\text{)} \\
&= \operatorname{Exp}(a\operatorname{Log}f + a\operatorname{Log}g) \\
&= \operatorname{Exp}(a\operatorname{Log}f)\cdot\operatorname{Exp}(a\operatorname{Log}g)
= f^a\cdot g^a.
\end{align*}

\textbf{Third rule ($(f^a)^b = f^{ab}$):}
\begin{align*}
(f^a)^b &= \operatorname{Exp}(b\operatorname{Log}(f^a)) \\
&= \operatorname{Exp}(b\operatorname{Log}(\operatorname{Exp}(a\operatorname{Log}f))) \\
&= \operatorname{Exp}(b\cdot a\operatorname{Log}f)
\quad\text{(by }\operatorname{Log}\circ\operatorname{Exp} = \operatorname{id}\text{)} \\
&= \operatorname{Exp}((ab)\operatorname{Log}f) = f^{ab}.
\end{align*}
\end{proof}

\begin{lemma}
\label{lem.fps.fpsPow-nat}
\lean{AlgebraicCombinatorics.FPS.fpsPow_nat}
\leanhelper
\leanok
For $f\in K[[x]]_1$ and $n\in\mathbb{N}$, we have
$f^n = \underbrace{f\cdot f\cdots f}_{n\text{ times}}$.
That is, our definition of $f^c$ is consistent with the standard power when $c\in\mathbb{N}$.
\end{lemma}

\begin{proof}
\leanok
By induction on $n$:
the base case $n=0$ gives $f^0=1$ (by Lemma~\ref{lem.fps.fpsPow-zero}),
and the inductive step uses $f^{n+1} = f^n\cdot f^1 = f^n\cdot f$
(by the first rule of exponents and Lemma~\ref{lem.fps.fpsPow-one}).
\end{proof}

\subsection{The Newton binomial formula for arbitrary exponents}

\begin{theorem}
[Generalized Newton binomial formula]
\label{thm.fps.gen-newton}
\lean{AlgebraicCombinatorics.FPS.generalizedNewtonBinomial}
\leantarget
\leanok
Assume that $K$
is a commutative $\mathbb{Q}$-algebra. Let $c\in K$. Then,%
\[
\left(  1+x\right)  ^{c}=\sum_{k\in\mathbb{N}}\dbinom{c}{k}x^{k}.
\]
\end{theorem}

\begin{proof}
The definition of $\operatorname{Log}$ yields
$\operatorname{Log}(1+x) = \overline{\log}\circ x = \overline{\log}$.
By Definition~\ref{def.fps.power-c}, we have
\begin{align*}
(1+x)^c &= \operatorname{Exp}(c\operatorname{Log}(1+x))
= \operatorname{Exp}(c\,\overline{\log}) \\
&= \exp\circ(c\,\overline{\log})
= \sum_{m\in\mathbb{N}} \frac{1}{m!}\left(\sum_{n\geq 1}\frac{(-1)^{n-1}}{n}cx^n\right)^m.
\end{align*}
Expanding and collecting by powers of $x$, we obtain
\[
(1+x)^c = \sum_{k\in\mathbb{N}} \left(\sum_{(n_1,\ldots,n_m)\in\operatorname{Comp}(k)}
\frac{1}{m!}\cdot\frac{(-1)^{n_1+\cdots+n_m-m}}{n_1 n_2\cdots n_m} c^m\right) x^k,
\]
where $\operatorname{Comp}(k)$ denotes the set of all compositions of~$k$.

It remains to show that the coefficient of $x^k$ equals $\binom{c}{k}$.
We use the \emph{polynomial identity trick}: for each $k\in\mathbb{N}$,
both sides of the coefficient equality
\[
\sum_{(n_1,\ldots,n_m)\in\operatorname{Comp}(k)}
\frac{1}{m!}\cdot\frac{(-1)^{n_1+\cdots+n_m-m}}{n_1\cdots n_m} c^m = \binom{c}{k}
\]
are polynomials in~$c$ with rational coefficients.
(The left-hand side is a finite sum of ``rational number times a power of~$c$''
expressions.
The right-hand side is $\frac{c(c-1)(c-2)\cdots(c-k+1)}{k!}$.)

For $c = n\in\mathbb{N}$, the equality holds by the standard Newton binomial
formula (Theorem~\ref{thm.fps.newton-binom}), since both sides equal $\binom{n}{k}$.
Since two polynomials with rational coefficients that agree on all natural
numbers must be equal (having infinitely many common values), the equality
holds for all $c\in K$.

In the formal proof:
\begin{itemize}
\item The base case $k=0$: both sides equal~$1$.
\item The inductive case $k=k'+1$: both the LHS and the RHS are expressed
as evaluations of explicit polynomials over~$\mathbb{Q}$, mapped to~$K$.
The LHS polynomial is built by expanding $\exp\circ(c\,\overline{\log})$
as $\sum_{d\ge 0}\frac{1}{d!}(c\,\overline{\log})^d$;
the scalar factoring $(c\,\overline{\log})^d = c^d\,\overline{\log}^d$
uses Lemma~\ref{lem.fps.smul-pow-eq-pow-smul},
and the truncation of the sum at $d=k$ uses the vanishing
$[x^k](\overline{\log}^d)=0$ for $d>k$
(Lemma~\ref{lem.fps.coeff-logSeries-pow-eq-zero}).
The polynomial identity principle (Lemma~\ref{lem.fps.poly-eq-of-nat-eval-eq}) then
shows these polynomials are equal, since they agree on all $n\in\mathbb{N}$.
\end{itemize}
\end{proof}

\begin{corollary}
\label{cor.fps.coeff-one-add-X-fpsPow}
\lean{AlgebraicCombinatorics.FPS.coeff_one_add_X_fpsPow}
\leanhelper
\leanok
For any $c\in K$ and $k\in\mathbb{N}$,
\[
[x^k](1+x)^c = \binom{c}{k}.
\]
\end{corollary}

\begin{proof}
\leanok
Immediate from Theorem~\ref{thm.fps.gen-newton} by reading off coefficients.
\end{proof}

\begin{theorem}
\label{thm.fps.antiNewtonBinomial}
\lean{AlgebraicCombinatorics.FPS.antiNewtonBinomial}
\leanhelper
\leanok
For any $n\in K$,
\[
(1+x)^{-n} = \sum_{i\in\mathbb{N}} (-1)^i \binom{n+i-1}{i} x^i.
\]
\end{theorem}

\begin{proof}
\leanok
Apply Theorem~\ref{thm.fps.gen-newton} with $c=-n$ and use the upper negation
identity $\binom{-n}{i} = (-1)^i\binom{n+i-1}{i}$.
\end{proof}

\subsection{Further rules of exponents}

\begin{lemma}
\label{lem.fps.one-fpsPow}
\lean{AlgebraicCombinatorics.FPS.one_fpsPow}
\leanhelper
\leanok
For any $c\in K$, we have $1^c = 1$.
\end{lemma}

\begin{proof}
\leanok
$1^c = \operatorname{Exp}(c\cdot\operatorname{Log}1) = \operatorname{Exp}(c\cdot 0) = \operatorname{Exp}(0) = 1$.
\end{proof}

\begin{lemma}
\label{lem.fps.constantCoeff-fpsPow}
\lean{AlgebraicCombinatorics.FPS.constantCoeff_fpsPow}
\leanhelper
\leanok
For $f\in K[[x]]_1$ and $c\in K$, we have $[x^0](f^c) = 1$.
\end{lemma}

\begin{proof}
\leanok
Immediate from $f^c\in K[[x]]_1$.
\end{proof}

\begin{lemma}
\label{lem.fps.fpsPow-neg}
\lean{AlgebraicCombinatorics.FPS.fpsPow_neg}
\leanhelper
\leanok
For $f\in K'[[x]]_1$ (over a field $K'$) and $c\in K'$,
\[
f^{-c} = (f^c)^{-1}.
\]
\end{lemma}

\begin{proof}
\leanok
We have $f^{-c}\cdot f^c = f^{-c+c} = f^0 = 1$, so $f^{-c}$ is the
inverse of~$f^c$.
\end{proof}

\begin{lemma}
\label{lem.fps.fpsPow-int}
\lean{AlgebraicCombinatorics.FPS.fpsPow_int}
\leanhelper
\leanok
For integer exponents, $f^c$ agrees with the standard power:
for $n\ge 0$, $f^n = f^n$ (standard), and for $n<0$, $f^n = (f^{|n|})^{-1}$.
\end{lemma}

\begin{proof}
\leanok
For $n\ge 0$, use Lemma~\ref{lem.fps.fpsPow-nat}.
For $n<0$, combine $f^{-c}\cdot f^c = 1$ with $f^{|n|}$ computed via the natural-number case.
\end{proof}

\subsection{Another application}

\begin{proposition}
\label{prop.binom.nCk-2i-qedmo.CN}
\lean{AlgebraicCombinatorics.FPS.binomialIdentity}
\leantarget
\leanok
Let $n\in\mathbb{C}$ and $k\in\mathbb{N}$.
Then,%
\[
\sum_{i=0}^{k}\dbinom{n+i-1}{i}\dbinom{n}{k-2i}=\dbinom{n+k-1}{k}.
\]
\end{proposition}

\begin{proof}
Define two FPSs $f,g\in\mathbb{C}\left[  \left[  x\right]  \right]  $ by
\begin{align*}
f  &  =\sum_{i\in\mathbb{N}}\dbinom{n+i-1}{i}x^{2i}\\
\text{and}\ \ \ \ \ \ \ \ \ \ g  &  =\sum_{j\in\mathbb{N}}\dbinom{n}{j}x^{j}.
\end{align*}
Multiplying and collecting by powers of $x$, we find that
\[
[x^k](fg) = \sum_{i=0}^{k} \dbinom{n+i-1}{i}\dbinom{n}{k-2i},
\]
which is the left-hand side of our identity.

For $g$, the generalized Newton formula (Theorem~\ref{thm.fps.gen-newton})
gives $g=(1+x)^n$.
For $f$, the anti-Newton binomial formula
(Theorem~\ref{thm.fps.antiNewtonBinomial}) gives $(1+x)^{-n} = \sum_{i\in\mathbb{N}}(-1)^i\binom{n+i-1}{i}x^i$.
Substituting $-x^2$ for $x$, we get $f = (1-x^2)^{-n}$.

Thus
\begin{align*}
fg &= (1-x^2)^{-n}(1+x)^n \\
&= \left(\frac{1+x}{1-x^2}\right)^n
= \left(\frac{1-x^2}{1+x}\right)^{-n} \\
&= (1-x)^{-n}.
\end{align*}
Substituting $-x$ for $x$ in the anti-Newton formula gives
\[
(1-x)^{-n} = \sum_{i\in\mathbb{N}}\binom{n+i-1}{i}x^i.
\]
Hence $[x^k](fg) = \binom{n+k-1}{k}$.

The formal proof establishes that $fg = (1-x)^{-n}$ by first showing $f = (1-x^2)^{-n}$.
The identity $f = (1-x^2)^{-n}$ is proved by comparing
coefficients using the polynomial identity principle:
for natural~$N$, both sides of the coefficient identity follow from
the inverse of $(1-x^2)^N$ expressed via expansion of $(1-x)^{-N}$, and uniqueness of inverses
gives agreement;
the polynomial identity principle then extends the result to all $n\in K$.
The odd-coefficient vanishing is established by a pairing argument: for odd~$k$,
each pair $(i,k-i)$ contributes terms with opposite signs (since $i$ and $k-i$ have
opposite parities), giving a sum of~$0$.
\end{proof}

\subsection{Helper lemmas from the formalization}

\begin{lemma}
\label{lem.fps.logSeries-eq-logbar}
\lean{AlgebraicCombinatorics.FPS.logSeries_eq_logbar}
\leanhelper
\leanok
Our $\overline{\log}$ series equals the logarithm series from
Mathlib's Exp/Log framework.
\end{lemma}

\begin{proof}
\leanok
Coefficient comparison.
\end{proof}

\begin{lemma}
\label{lem.fps.fpsLog-eq-Log-val}
\lean{AlgebraicCombinatorics.FPS.fpsLog_eq_Log_val}
\leanhelper
\leanok
The definition of $\operatorname{Log}$ from Definition~\ref{def.fps.fpsLog}
is equivalent to the definition from Mathlib's Exp/Log framework for FPS with constant term one.
\end{lemma}

\begin{proof}
\leanok
Unfold definitions and use Lemma~\ref{lem.fps.logSeries-eq-logbar}.
\end{proof}

\begin{lemma}
\label{lem.fps.fpsExp-eq-Exp-val}
\lean{AlgebraicCombinatorics.FPS.fpsExp_eq_Exp_val}
\leanhelper
\leanok
For $g$ with $[x^0]g=0$, the definition of $\operatorname{Exp}(g)$ from Definition~\ref{def.fps.fpsExp}
is equivalent to the definition from Mathlib's Exp/Log framework for FPS with constant term zero.
\end{lemma}

\begin{proof}
\leanok
Unfold definitions.
\end{proof}

\begin{lemma}
\label{lem.fps.smul-pow-eq-pow-smul}
\lean{AlgebraicCombinatorics.FPS.smul_pow_eq_pow_smul}
\leanhelper
\leanok
$(c\cdot g)^m = c^m\cdot g^m$ for $c\in K$, $g\in K[[x]]$, $m\in\mathbb{N}$.
\end{lemma}

\begin{proof}
\leanok
Induction on $m$.
\end{proof}

\begin{lemma}
\label{lem.fps.coeff-logSeries-pow-eq-zero}
\lean{AlgebraicCombinatorics.FPS.coeff_logSeries_pow_eq_zero_of_gt}
\leanhelper
\leanok
For $k < m$, we have $[x^k](\overline{\log}^m) = 0$
(since $\overline{\log}$ has order~$\geq 1$, so $\overline{\log}^m$
has order~$\geq m$).
\end{lemma}

\begin{proof}
\leanok
The order of $\overline{\log}^m$ is at least $m$ since
$[x^0]\overline{\log}=0$.
\end{proof}

\begin{theorem}
\label{thm.fps.polynomial-identity-trick}
\lean{AlgebraicCombinatorics.FPS.polynomial_identity_trick}
\leanhelper
\leanok
Let $R$ be a commutative domain of characteristic~$0$,
and let $p\in R[x]$ be a polynomial such that $p(n)=0$ for all $n\in\mathbb{N}$.
Then $p=0$.
\end{theorem}

\begin{proof}
\leanok
The set of roots of $p$ contains all natural numbers, which is an infinite subset
of~$R$ (since $R$ has characteristic~$0$). A nonzero polynomial of degree~$d$
has at most~$d$ roots, so $p$ must be zero.
\end{proof}

\begin{lemma}
\label{lem.fps.poly-eq-of-nat-eval-eq}
\lean{AlgebraicCombinatorics.FPS.poly_eq_of_nat_eval_eq}
\leanhelper
\leanok
If two polynomials over~$\mathbb{Q}$ agree when evaluated at all natural numbers
(after mapping to~$K$ via the algebra map), then they are equal:
if $p(n) = q(n)$ for all $n\in\mathbb{N}$,
then $p = q$.
\end{lemma}

\begin{proof}
Apply the polynomial identity trick to $p - q$.
\end{proof}

\begin{lemma}
\label{lem.fps.Ring-choose-pascal}
\lean{AlgebraicCombinatorics.FPS.Ring_choose_pascal}
\leanhelper
\leanok
Pascal's rule for generalized binomial coefficients:
\[
\binom{c}{k+1} = \binom{c-1}{k} + \binom{c-1}{k+1}.
\]
\end{lemma}

\begin{proof}
Follows from the Pascal recurrence for generalized binomial coefficients.
\end{proof}

\begin{lemma}
\label{lem.fps.fpsPow-coeff-pascal}
\lean{AlgebraicCombinatorics.FPS.fpsPow_coeff_pascal}
\leanhelper
\leanok
The coefficients of $(1+x)^c$ satisfy Pascal's recurrence:
\[
[x^{k+1}](1+x)^c = [x^k](1+x)^{c-1} + [x^{k+1}](1+x)^{c-1}.
\]
\end{lemma}

\begin{proof}
Since $(1+x)^{c} = (1+x)^{c-1}\cdot(1+x)$, the coefficient $[x^{k+1}]$ of the
product is $[x^k](1+x)^{c-1} + [x^{k+1}](1+x)^{c-1}$.
\end{proof}

\subsection{Helpers for Theorem~\ref{thm.fps.gen-newton}}

The generalized Newton binomial formula requires showing that both the LHS
(coefficients of $(1+x)^c$) and the RHS ($\binom{c}{k}$) are polynomial
functions of~$c$ that agree on all natural numbers.

\begin{lemma}
\label{lem.fps.Ring-choose-eq-choosePoly-eval}
\lean{AlgebraicCombinatorics.FPS.Ring_choose_eq_choosePoly_eval}
\leanhelper
\leanok
The generalized binomial coefficient $\binom{c}{k}$ equals the evaluation
of a polynomial in~$c$: there exists a polynomial $p_k\in\mathbb{Q}[x]$ such that
$\binom{c}{k} = p_k(c)$ for all $c\in K$.
\end{lemma}

\begin{proof}
The polynomial is $p_k(x) = x(x-1)\cdots(x-k+1)/k!$,
evaluated at~$c$ via the algebra map $\mathbb{Q} \to K$.
\end{proof}

\begin{lemma}
\label{lem.fps.coeff-fpsPow-eq-coeffExpPoly-eval}
\lean{AlgebraicCombinatorics.FPS.coeff_fpsPow_eq_coeffExpPoly_eval}
\leanhelper
\leanok
The coefficient $[x^k](1+x)^c$ equals the evaluation of a polynomial in~$c$:
\[
[x^k](1+x)^c = q_k(c)
\]
for some $q_k\in\mathbb{Q}[x]$.
\end{lemma}

\begin{proof}
Expand $(1+x)^c = \operatorname{Exp}(c\cdot\overline{\log})
= \sum_{d\ge 0}\frac{1}{d!}(c\,\overline{\log})^d$.
Using $(c\,\overline{\log})^d = c^d\,\overline{\log}^d$
and the vanishing $[x^k](\overline{\log}^d)=0$ for $d>k$,
the coefficient $[x^k]$ is a polynomial in~$c$ of degree~$\leq k$.
\end{proof}

\begin{lemma}
\label{lem.fps.one-sub-X-pow-neg-coeff}
\lean{AlgebraicCombinatorics.FPS.one_sub_X_pow_neg_coeff}
\leanhelper
\leanok
For $n\in K$ and $k\in\mathbb{N}$,
\[
[x^k](1-x)^{-n} = \binom{n+k-1}{k}.
\]
\end{lemma}

\begin{proof}
\leanok
Use the upper negation identity $\binom{-n}{k} = (-1)^k\binom{n+k-1}{k}$
and the sign cancellation $(-1)^k\cdot(-1)^k = 1$.
\end{proof}

\subsection{Helpers for Proposition~\ref{prop.binom.nCk-2i-qedmo.CN}}

The proof of the binomial identity requires establishing that the product
$f\cdot g$ (where $f = (1-x^2)^{-n}$ and $g = (1+x)^n$) equals $(1-x)^{-n}$.
This is done in several steps:

\begin{enumerate}
\item Define $f = \sum_{i\in\mathbb{N}} \binom{n+i-1}{i} x^{2i}$ and
      $g = \sum_{j\in\mathbb{N}} \binom{n}{j} x^j$.
\item Show $fg = (1-x)^{-n}$
      (Lemma~\ref{lem.fps.key-product-identity}).
\item Read off coefficients.
\end{enumerate}

\begin{definition}
\label{def.fps.f-series-g-series}
\lean{AlgebraicCombinatorics.FPS.f_series', AlgebraicCombinatorics.FPS.g_series'}
\leanhelper
\leanok
Define two FPSs for the binomial identity:
\begin{align*}
f(n) &:= \sum_{i\in\mathbb{N}} \binom{n+i-1}{i} x^{2i}
\quad\text{(i.e., the even-indexed coefficients of $(1-x^2)^{-n}$)}, \\
g(n) &:= \sum_{j\in\mathbb{N}} \binom{n}{j} x^j = (1+x)^n.
\end{align*}
\end{definition}

\begin{lemma}
\label{lem.fps.coeff-f-mul-g}
\lean{AlgebraicCombinatorics.FPS.coeff_f_mul_g'}
\leanhelper
\leanok
The coefficient of $x^k$ in $f(n)\cdot g(n)$ equals
\[
[x^k](f\cdot g) = \sum_{i=0}^{\lfloor k/2\rfloor} \binom{n+i-1}{i}\binom{n}{k-2i}.
\]
\end{lemma}

\begin{proof}
By expanding the Cauchy product and using the fact that $f$ has
nonzero coefficients only at even positions.
\end{proof}

\begin{lemma}
\label{lem.fps.invOneSubPow-mul-rescale-inv}
\lean{AlgebraicCombinatorics.FPS.invOneSubPow_mul_rescale_mul_one_sub_sq_pow'}
\leanhelper
\leanok
For natural $N$,
\[
(1-x)^{-N}\cdot (1+x)^{-N}\cdot (1-x^2)^N = 1.
\]
That is, $(1-x)^{-N}\cdot(1+x)^{-N}$ is the inverse of $(1-x^2)^N$.
\end{lemma}

\begin{proof}
Factor $(1-x^2)^N = (1-x)^N(1+x)^N$ and use the inverse relations
$(1-x)^{-N}\cdot(1-x)^N = 1$ and $(1+x)^{-N}\cdot(1+x)^N = 1$.
\end{proof}

\begin{lemma}
\label{lem.fps.expand-invOneSubPow-inv}
\lean{AlgebraicCombinatorics.FPS.expand_invOneSubPow_eq_inv_one_sub_X_sq}
\leanhelper
\leanok
For natural $N$,
$(1-x)^{-N}\big|_{x \mapsto x^2}\cdot (1-x^2)^N = 1$.
That is, $(1-x)^{-N}\big|_{x \mapsto x^2}$ is also the inverse of $(1-x^2)^N$.
\end{lemma}

\begin{proof}
Since the substitution $x \mapsto x^2$ is a ring homomorphism and $(1-x)\big|_{x \mapsto x^2} = 1-x^2$,
we have $(1-x)^{-N}\big|_{x \mapsto x^2}\cdot(1-x)^N\big|_{x \mapsto x^2}
= 1\big|_{x \mapsto x^2} = 1$.
\end{proof}

\begin{lemma}
\label{lem.fps.product-eq-expand}
\lean{AlgebraicCombinatorics.FPS.product_eq_expand_invOneSubPow}
\leanhelper
\leanok
For natural $N$,
\[
(1-x)^{-N}\cdot(1+x)^{-N}
= (1-x)^{-N}\big|_{x \mapsto x^2}.
\]
Both sides are inverses of $(1-x^2)^N$, so they are equal by uniqueness of inverses.
\end{lemma}

\begin{proof}
Both $P := (1-x)^{-N}\cdot(1+x)^{-N}$ and
$E := (1-x)^{-N}\big|_{x \mapsto x^2}$ satisfy $(\cdot)\cdot(1-x^2)^N = 1$.
Since $(1-x^2)^N$ is a unit, $P = E$.
\end{proof}

\begin{lemma}
\label{lem.fps.coeff-invOneSubPow-product-even}
\lean{AlgebraicCombinatorics.FPS.coeff_invOneSubPow_product_even'}
\leanhelper
\leanok
For natural $N\ge 1$ and $m\in\mathbb{N}$,
\[
[x^{2m}]\bigl((1-x)^{-N}\cdot(1+x)^{-N}\bigr) = \binom{N+m-1}{m}.
\]
\end{lemma}

\begin{proof}
By Lemma~\ref{lem.fps.product-eq-expand}, the product equals
$(1-x)^{-N}\big|_{x \mapsto x^2}$, so the coefficient at $2m$ is
$[x^m](1-x)^{-N} = \binom{N+m-1}{m}$.
\end{proof}

\begin{lemma}[Key product identity]
\label{lem.fps.key-product-identity}
\lean{AlgebraicCombinatorics.FPS.key_product_identity'}
\leanhelper
\leanok
For all $n\in K$,
\[
f(n)\cdot g(n) = (1-x)^{-n}.
\]
In other words, $(1-x^2)^{-n}\cdot(1+x)^n = (1-x)^{-n}$.
\end{lemma}

\begin{proof}
First show $f(n) = (1-x^2)^{-n}$ by establishing that $f(n)$ equals
$(1+x)^{-n} \cdot (1-x)^{-n}$.
The even-coefficient identity uses Lemma~\ref{lem.fps.coeff-invOneSubPow-product-even} for natural~$N$
and then extends to all~$n$ by the polynomial identity principle
(Theorem~\ref{thm.fps.polynomial-identity-trick}).
The odd-coefficient vanishing follows from a pairing argument: for odd~$k$,
each pair $(i,k-i)$ contributes terms with opposite signs.
Then $f(n)\cdot g(n) = (1-x^2)^{-n}\cdot(1+x)^n$, and the cancellation
$(1+x)^{-n}\cdot(1+x)^n = 1$
gives the result.
\end{proof}
